\newtheorem{definition}{Definition}
\newtheorem{example}{Example}

\newtheorem{proposition}{Proposition}


\newcommand{\Reals}{\mathbb{R}}
\newcommand{\Nats}{\mathbb{N}}
\newcommand{\bvec}[1]{\boldsymbol{#1}}
\newcommand{\Dset}{\mathcal{D}}
\newcommand{\Pol}{\mathbf{P}}

\newcommand{\Argmax}{\text{argmax}}
\newcommand{\PositiveReals}{\Reals_{\geq 0}}

\newcommand{\Mass}[1]{\Sigma #1}
 \newcommand{\defsymbol}{\stackrel{\text{def}}{=}}

\newcommand{\support}{\boldsymbol{s}}

% This wrapper is just to avoid some wrong warnings in overleaf about math mode
\def\oo{\infty}
\def\eps{\epsilon}
\def\calX{{\mathcal{X}}}

%\NewDocumentCommand{\MEC}{O{\alpha,\beta,\pi}}{\text{MEC}_{#1}^\pi}

\NewDocumentCommand{\ER}{O{\alpha}}{\operatorname{ER}_{#1}}

% Candidate polarization measure labels used in prose and compact comparisons.
\NewDocumentCommand{\ERLabel}{o}{\ensuremath{\mathrm{ER}\IfValueT{#1}{_{#1}}}}
\newcommand{\EMDLabel}{\ensuremath{\mathrm{EMD}}}
\newcommand{\VdELabel}{\ensuremath{\mathrm{VdE}}}
\newcommand{\TWDLabel}{\ensuremath{\mathrm{TWD}}}
\newcommand{\KILabel}{\ensuremath{\mathrm{KI}}}

\NewDocumentCommand{\EC}{O{\alpha,\beta} O{\pi}}{%
\operatorname{EC}_{#1}^{#2}%
}

\NewDocumentCommand{\MEC}{O{\alpha,\beta}}{\operatorname{MEC}_{#1}}

\NewDocumentCommand{\ECEMD}{O{\alpha,\beta} O{\pi}}{%
\operatorname{EC}^{#2}%
}

\NewDocumentCommand{\MECEMD}{O{\alpha,\beta}}{\operatorname{MEC}}

\def\EMD{\operatorname{EMD}}

\NewDocumentCommand{\set}{s m}{\IfBooleanTF{#1}{{\{#2\}}}{{\ensuremath{\left\{#2\right\}}}}}
\NewDocumentCommand{\lr}{s m}{\IfBooleanTF{#1}{{(#2)}}{{\ensuremath{\left(#2\right)}}}}
\NewDocumentCommand{\eqdef}{}{\stackrel{{\scriptscriptstyle \mathrm{def}}}{=}}
\NewDocumentCommand{\argmin}{}{\operatorname*{arg\,min}}
\NewDocumentCommand{\bpi}{}{{\boldsymbol{\pi}}}
\NewDocumentCommand{\bx}{}{\boldsymbol{x}}
\NewDocumentCommand{\midrange}{}{{\sigma_\oo}}
\NewDocumentCommand{\lrfrac}{m m}{\lr{\frac{#1}{#2}}}


\newcommand{\doubleCheckmark}[2][1]{%
  \scalebox{#1}{\checkmark\hspace{-#2}\textcolor{white}{\checkmark}\hspace{-#2}\checkmark}%
}
\def\checkcheck{{\doubleCheckmark{1em}}}

\newcommand{\Var}{\mathrm{Var}}

\newcommand{\Disc}{\Delta
}
% \NewDocumentCommand{\pixn}{O{}}{(\pi^{#1}_i)_{i=1}^n, (x_i)_{i=1}^n}

\newcommand{\proofpart}[1]{\par\medskip\noindent\textbf{#1.}\ }
